%!TEX root = main.tex
% ==============================================================================
% İSTANBUL ÜNİVERSİTESİ SOSYAL BİLİMLER ENSTİTÜSÜ
% BÖLÜM 2 (Madde 16/c, d ve Madde 17/b)
% ==============================================================================
\chapter{VERİ ANALİZİ VE BULGULAR}

\section{Araştırma Bulgularının Sınıflandırılması}

Bu bölümde toplanan ampirik veriler ve kuramsal karşılaştırmalar yer almaktadır.\footnote{Yeni bölüme geçildiğinde dipnot numaralandırması Madde 19/a gereğince tekrar 1'den başlar.} Sosyolojik sistem analizlerinde sıklıkla başvurulan yaklaşımlardan biri Parsons tarafından geliştirilmiştir \parencite{parsons_social_system}.

\begin{table}[htbp]
    \centering
    \caption{Araştırmada İncelenen Değişkenlerin Dağılımı}
    \label{tab:degiskenler}
    \begin{tabular}{|l|c|r|}
        \hline
        \textbf{Kategori} & \textbf{Örneklem Sayısı (N)} & \textbf{Oran (\%)} \\
        \hline
        Birinci Grup & 120 & 40,0 \\
        İkinci Grup  & 150 & 50,0 \\
        Üçüncü Grup  & 30  & 10,0 \\
        \hline
        \textbf{Toplam} & \textbf{300} & \textbf{100,0} \\
        \hline
    \end{tabular}
\end{table}

Tablo \ref{tab:degiskenler} incelendiğinde, örneklem gruplarının oransal dağılımı görülmektedir. Tablo ve şekillerin numaralandırılmasında Arap veya Romen rakamları tutarlı biçimde kullanılır (Madde 3/b).

\section{Kurumsal Yapı ve Karşılaştırmalı Değerlendirme}

Türkiye'de anayasal gelişmeler ve insan hakları kuramları incelenirken yerli kaynaklara ve süreli yayınlara başvurulmuştur \parencite{tanor_osmanli, kaboglu_dusunce}.

\begin{figure}[htbp]
    \centering
    \framebox[0.75\textwidth][c]{\rule{0pt}{4cm} Şekil / Grafik Alanı}
    \caption{Araştırma Modeli ve Değişkenler Arası Etkileşim}
    \label{fig:model}
\end{figure}

Şekil \ref{fig:model}'de görüldüğü üzere modelin kavramsal boyutları birbiriyle bütüncül bir etkileşim içerisindedir. Şekil ve tablo listeleri ön kısımda tek satır aralığı ile yer almaktadır (Madde 4/b).